\documentclass[11pt,a4paper]{article}

\usepackage{amsmath, amssymb}
\usepackage{geometry}
\usepackage{hyperref}
\usepackage{booktabs}
\usepackage{array}
\usepackage{xcolor}
\usepackage{parskip}
\usepackage{titlesec}
\usepackage{fancyhdr}
\usepackage{microtype}
\usepackage{enumitem}
\usepackage{listings}

\geometry{margin=2.4cm}

\definecolor{codebg}{HTML}{F5F5F2}
\definecolor{rulec}{HTML}{C3C2B7}
\definecolor{warn}{HTML}{9C4221}

\hypersetup{colorlinks=true, linkcolor=blue!55!black, urlcolor=blue!55!black}

\pagestyle{fancy}
\fancyhf{}
\rhead{\small AutoDock Vina --- Practical Tutorial}
\lhead{\small \nouppercase{\leftmark}}
\cfoot{\thepage}

\titleformat{\section}{\large\bfseries}{\thesection}{0.8em}{}[\vspace{-5pt}\rule{\linewidth}{0.4pt}]
\titleformat{\subsection}{\normalsize\bfseries}{\thesubsection}{0.8em}{}

\lstset{
  basicstyle=\ttfamily\small,
  backgroundcolor=\color{codebg},
  frame=single,
  rulecolor=\color{rulec},
  breaklines=true,
  columns=fullflexible,
  xleftmargin=0pt,
  framexleftmargin=6pt,
  framexrightmargin=6pt,
  aboveskip=10pt,
  belowskip=10pt,
}

\newcommand{\code}[1]{\texttt{#1}}
\newcommand{\pitfall}[1]{\textcolor{warn}{\textbf{#1}}}

\title{\vspace{-1.5cm}\textbf{AutoDock Vina: A Practical Tutorial}\\[6pt]
  \large Preparing inputs, running docking, and checking that the answer is real}
\author{CADDack Project}
\date{\today}

\begin{document}
\maketitle
\thispagestyle{fancy}

\begin{abstract}
\noindent
A working guide to molecular docking with AutoDock Vina: what the program actually
computes, how to turn a PDB and a SMILES string into valid inputs, how to read the
output, and --- the part most tutorials skip --- how to verify that your setup
produces correct poses rather than merely running without errors. Code examples use
the \code{caddack.docking} module, but every step maps directly onto plain Vina.
The pitfalls in Section~\ref{sec:pitfalls} are ones measured in practice, with the
numbers they cost.
\end{abstract}

% =====================================================================
\section{What docking is, and what Vina computes}
% =====================================================================

Docking answers a geometric question: \emph{given a protein structure and a small
molecule, where and in what orientation does the molecule bind?} A docking program
has two parts:

\begin{enumerate}[leftmargin=1.3em, itemsep=2pt]
  \item a \textbf{search} over the ligand's position, orientation and flexible
        torsions, and
  \item a \textbf{scoring function} that ranks candidate poses.
\end{enumerate}

Vina's scoring function is empirical and pairwise. For every pair of heavy atoms
within a cutoff it sums terms depending only on the surface distance
$d_{ij} = r_{ij} - R_i - R_j$:

\begin{center}
\renewcommand{\arraystretch}{1.25}
\begin{tabular}{@{}ll@{}}
\toprule
\textbf{Term} & \textbf{Role} \\
\midrule
\code{gauss1}, \code{gauss2} & steric attraction --- rewards shape complementarity \\
\code{repulsion} & penalises overlap ($d_{ij} < 0$) \\
\code{hydrophobic} & rewards apolar carbon contact \\
\code{hydrogen bonding} & rewards donor--acceptor pairs at short range \\
\code{$N_{\text{rot}}$ penalty} & entropic cost of freezing rotatable bonds \\
\bottomrule
\end{tabular}
\end{center}

The reported affinity is in kcal/mol; more negative means predicted stronger
binding.

\subsection*{What Vina does \emph{not} use}

Two omissions matter for how you prepare inputs.

\pitfall{Partial charges are ignored.} The Vina scoring function has no
electrostatic term --- everything above depends on \emph{atom types} and distances.
The charge column in a PDBQT file is read but unused. Charges only matter if you
switch to the AutoDock4 scoring function (\code{sf\_name="ad4"}), which does have an
electrostatic term.

\pitfall{Ring conformations are never changed.} Vina samples translation, rotation
and \emph{acyclic} torsions only. Whatever ring pucker your input conformer has is
what you get. Section~\ref{sec:pitfalls} shows what this costs when ignored.

% =====================================================================
\section{The PDBQT format}
% =====================================================================

Vina does not read PDB. It reads \textbf{PDBQT}: a PDB with two extra pieces of
information --- a partial charge (Q) and an AutoDock atom type (T).

\begin{lstlisting}
ATOM      1  N   TYR A  20      22.423  16.962  23.645  1.00  0.00    -0.346 N
|<-1-6->|<7-11>| |13-16|17|18-20|21|22|23-26|  31-38   55-60 61-66     71-76  78-79
\end{lstlisting}

The layout is column-strict --- a field off by one character makes Vina reject the
file with a parsing error. The two added columns are charge (71--76) and atom type
(78--79).

\subsection*{Atom types}

Typing is what actually drives the scoring function, so getting it right matters
more than the charges do.

\begin{center}
\renewcommand{\arraystretch}{1.2}
\begin{tabular}{@{}llp{7.6cm}@{}}
\toprule
\textbf{Type} & \textbf{Element} & \textbf{Meaning} \\
\midrule
\code{C}  & carbon   & aliphatic, hydrophobic \\
\code{A}  & carbon   & aromatic \\
\code{N}  & nitrogen & \emph{not} a hydrogen-bond acceptor \\
\code{NA} & nitrogen & hydrogen-bond \emph{acceptor} \\
\code{OA} & oxygen   & hydrogen-bond acceptor (almost all oxygens) \\
\code{SA} & sulfur   & acceptor \\
\code{HD} & hydrogen & polar hydrogen --- marks a hydrogen-bond \emph{donor} \\
\code{Zn}, \code{Mg}, \dots & metals & coordinating centres \\
\bottomrule
\end{tabular}
\end{center}

A donor is represented by its attached \code{HD} hydrogen, not by the heavy atom.
This is why a receptor stripped of hydrogens cannot donate hydrogen bonds at all:
if you remove every H and mark every nitrogen \code{NA}, you have converted every
backbone amide from a donor into an acceptor and inverted much of the site's
hydrogen-bonding character.

% =====================================================================
\section{Preparing the ligand}
% =====================================================================

The ligand needs 3D coordinates, hydrogens, correct bond orders, and a torsion
tree marking which bonds Vina may rotate.

\begin{lstlisting}[language=Python]
from caddack.docking import ligand_pdbqt_from_smiles

ligand = ligand_pdbqt_from_smiles(
    "CC(=O)Oc1ccccc1C(=O)O",     # aspirin
    "ligand.pdbqt",
)
\end{lstlisting}

Internally this is three steps:

\begin{enumerate}[leftmargin=1.3em, itemsep=2pt]
  \item \textbf{Parse and add hydrogens.} Geometry depends on them, and polar
        hydrogens become the \code{HD} donor markers.
  \item \textbf{Embed a 3D conformer} (RDKit ETKDG) and relax it with MMFF94, or
        UFF for atom types MMFF does not parameterise.
  \item \textbf{Write PDBQT} (Meeko), which assigns Gasteiger charges, decides atom
        types, and builds the torsion tree.
\end{enumerate}

The torsion tree appears in the file as \code{ROOT}/\code{ENDROOT} around the rigid
core with \code{BRANCH}/\code{ENDBRANCH} for each rotatable bond. \code{TORSDOF} at
the end gives the count used in the entropic penalty.

\subsection*{Protonation is a modelling decision}

A SMILES string does not fix the protonation state. \code{CC(=O)O} is neutral acetic
acid; \code{CC(=O)[O-]} is acetate. They have different atom types, different
hydrogen-bonding behaviour, and dock differently. Decide the state at your target pH
before docking --- tools such as \code{Dimorphite-DL} or Epik enumerate plausible
states.

% =====================================================================
\section{Preparing the receptor}
% =====================================================================

This is the fragile step, because a crystal structure is missing most of what Vina
needs: \textbf{no hydrogens, no bond orders, no charges, no atom types.}

\begin{lstlisting}[language=Python]
from caddack.docking import receptor_pdbqt_from_pdb

receptor = receptor_pdbqt_from_pdb(
    "protein.pdb", "receptor.pdbqt",
    keep_hetatm=["ZN"],     # keep a catalytic zinc; drop everything else
)
\end{lstlisting}

\subsection*{Cleaning}

Remove waters and crystallisation additives (glycerol, sulfate, PEG, DMSO), drop
alternate conformations, and keep a single model from an NMR ensemble. Keep
cofactors and metals that are part of the binding site --- deleting a catalytic zinc
removes the strongest anchor in the pocket.

Structural waters that bridge ligand and protein are a genuine judgement call:
Vina treats the receptor as rigid and dry, so a retained water is a fixed obstacle,
while a deleted one leaves a cavity that may be filled incorrectly.

\subsection*{Typing and charges come from residue templates}

Because a PDB gives no bond orders, charges cannot be computed from topology the way
they are for a ligand. Instead each residue is matched against a \textbf{template
library} carrying pre-assigned charges and types (this is what Meeko's
\code{ResidueChemTemplates} provides). Consequences:

\begin{itemize}[leftmargin=1.3em, itemsep=2pt]
  \item Non-standard residues, covalent modifications and interrupted chains fail to
        match, and there is no fallback that is remotely as good.
  \item A \emph{truncated pocket file} is not a valid polymer --- template matching
        needs whole residues in a connected chain. Prepare the complete protein.
\end{itemize}

\subsection*{Histidine, and why it needs a decision}

Histidine has three plausible states at physiological pH, and a crystal structure
cannot distinguish them because it has no hydrogens:

\begin{center}
\begin{tabular}{@{}lll@{}}
\toprule
\textbf{Template} & \textbf{Protonated at} & \textbf{Net charge} \\
\midrule
\code{HIE} & N$\varepsilon$2 & neutral (common default) \\
\code{HID} & N$\delta$1 & neutral \\
\code{HIP} & both & $+1$ \\
\bottomrule
\end{tabular}
\end{center}

The choice flips which ring nitrogen is a donor and which is an acceptor, so it
directly changes docking. Defaulting every histidine to \code{HIE} is a reasonable
convention; a histidine coordinating a metal or in a catalytic triad may well need
\code{HID} or \code{HIP}. For a principled assignment use \textbf{PDB2PQR with
PROPKA} (predicts per-residue p$K_\text{a}$) or \textbf{Reduce} (adds hydrogens and
optimises the hydrogen-bond network, including Asn/Gln/His flips).

% =====================================================================
\section{Defining the search box}
% =====================================================================

Vina searches only inside an axis-aligned box, given as a centre and a size in
\aa ngstr\"oms.

\begin{lstlisting}[language=Python]
from caddack.docking import box_from_reference_ligand, box_from_pocket_pdb

box = box_from_reference_ligand("known_ligand.mol2", padding=4.0)
box = box_from_pocket_pdb("pocket.pdb", padding=2.0)
\end{lstlisting}

\begin{itemize}[leftmargin=1.3em, itemsep=2pt]
  \item \textbf{Known site (re-docking, lead optimisation).} Centre on a
        co-crystallised ligand and pad by roughly 4\,\aa. Tight boxes search well.
  \item \textbf{Known pocket, no ligand.} Use the pocket residues' bounding box, or
        a cavity detector such as fpocket.
  \item \textbf{Unknown site (blind docking).} A box covering the whole protein
        works but searches far more space --- raise \code{exhaustiveness}
        substantially and expect worse ranking.
\end{itemize}

The box must comfortably contain the ligand in any orientation: a rule of thumb is
the ligand's longest dimension plus about 8\,\aa. Vina warns if the box is very
large, because search difficulty grows with volume.

% =====================================================================
\section{Running the docking}
% =====================================================================

\begin{lstlisting}[language=Python]
from caddack.docking import dock_pdbqt, dock_smiles

poses = dock_pdbqt(receptor, ligand, box,
                   exhaustiveness=32, n_poses=9, seed=42)

# or end-to-end from a SMILES string
poses = dock_smiles("protein.pdb", "CC(=O)Oc1ccccc1C(=O)O",
                    reference_ligand="known_ligand.mol2")

for p in poses:
    print(p.rank, p.score)      # kcal/mol, best first
\end{lstlisting}

\textbf{\code{exhaustiveness}} controls search effort. Vina's default is 8; 32 is
meaningfully more reliable for roughly four times the runtime, and blind docking
needs more still. It is the first parameter to raise when results look unstable.

\textbf{Determinism.} Vina is stochastic. A fixed \code{seed} makes single-threaded
runs reproducible, but multi-threaded runs may still vary slightly between
executions --- so treat a single docking run as one sample, not a measurement.

% =====================================================================
\section{Reading the output}
% =====================================================================

Each pose carries a predicted affinity and two RMSD values relative to the
best-scoring pose (a spread measure, \emph{not} accuracy against a crystal
structure).

\begin{center}
\renewcommand{\arraystretch}{1.2}
\begin{tabular}{@{}rrl@{}}
\toprule
\textbf{Rank} & \textbf{Affinity (kcal/mol)} & \textbf{Interpretation} \\
\midrule
1 & $-10.1$ & best-scoring pose \\
2 & $-9.4$  & plausible alternative \\
3 & $-7.2$  & clearly worse \\
\bottomrule
\end{tabular}
\end{center}

Two cautions on the number itself:

\begin{itemize}[leftmargin=1.3em, itemsep=2pt]
  \item \textbf{It is not a measured binding constant.} Vina's scoring function
        correlates only modestly with experiment --- roughly $r \approx 0.6$ on
        standard benchmarks. Use it to rank poses of \emph{one} ligand confidently,
        and to rank \emph{different} ligands only cautiously.
  \item \textbf{Scores drift with ligand size.} Larger molecules make more contacts
        and score better almost mechanically. Ligand-efficiency measures
        (score per heavy atom) partly correct for this.
\end{itemize}

A useful diagnostic is the score gap: if ranks 1 and 2 differ by only
$0.2$\,kcal/mol but sit in different orientations, the ranking is not meaningful and
both should be considered.

% =====================================================================
\section{Validating your setup: re-docking}
\label{sec:redock}
% =====================================================================

Before trusting docking on a new molecule, check that it reproduces answers you
already know. \textbf{Re-docking} takes a complex whose crystal structure is solved,
removes the ligand, docks it back, and measures RMSD to the crystal pose. Below
2\,\aa\ is the conventional success criterion.

\begin{lstlisting}[language=Python]
from caddack.docking import pose_rmsd
from rdkit import Chem

reference = Chem.MolFromMol2File("crystal_ligand.mol2")
print(pose_rmsd(poses[0], reference))    # symmetry-aware, in angstroms
\end{lstlisting}

RMSD must be \textbf{symmetry-aware}: a symmetric ring flipped 180$^\circ$ is the
same physical pose, and naive atom-order RMSD would score it as a failure. RDKit's
\code{CalcRMS} matches atoms by substructure and handles this.

\subsection*{Start from the crystal conformer}

Re-docking conventionally uses the deposited ligand conformer as the input. This is
not circular: Vina randomises position, orientation and acyclic torsions regardless.
What it \emph{does} preserve is the ring conformation --- and since Vina cannot
change ring puckers, re-generating the conformer instead measures conformer
generation and docking together, which is a different and harder question.

\subsection*{Reference numbers}

Measured on 30 randomly sampled PDBbind refined-set complexes, all of which docked
without error:

\begin{center}
\renewcommand{\arraystretch}{1.25}
\begin{tabular}{@{}lccc@{}}
\toprule
\textbf{Protocol} & \textbf{Top pose $<2$\,\aa} & \textbf{Any of 9 $<2$\,\aa} & \textbf{Median RMSD} \\
\midrule
Re-embedded ligand, \code{exh}\,$=8$   & 43\% & 70\% & 2.82\,\aa \\
Crystal conformer, \code{exh}\,$=8$    & 53\% & 83\% & 1.50\,\aa \\
Crystal conformer, \code{exh}\,$=32$   & 60\% & 90\% & 1.47\,\aa \\
\bottomrule
\end{tabular}
\end{center}

The gap between the last two columns is informative: the search usually
\emph{generates} a near-native pose, and the scoring function fails to rank it
first about a third of the time. That is the standard motivation for re-scoring a
pose set with a better model.

% =====================================================================
\section{Common pitfalls}
\label{sec:pitfalls}
% =====================================================================

\begin{enumerate}[leftmargin=1.4em, itemsep=6pt]

\item \pitfall{Re-generating the ligand conformer and expecting Vina to fix the
rings.} It cannot. In the table above this alone cost 10 percentage points. The
diagnostic signature is a ligand in the \emph{right place} with the wrong shape ---
one sugar sat $0.67$\,\aa\ from the correct position yet scored $4.37$\,\aa\ RMSD;
a macrocycle went from $7.67$\,\aa\ to $0.38$\,\aa\ once started from the deposited
conformer. If your ligand has flexible or macrocyclic rings, dock several ring
conformers.

\item \pitfall{Leaving \code{exhaustiveness} at 8 for anything demanding.} Worth
about 7 percentage points here for $4\times$ the runtime.

\item \pitfall{Preparing a receptor with no polar hydrogens.} If every nitrogen is
typed \code{NA} and there are no \code{HD} atoms, the receptor cannot donate a
hydrogen bond anywhere. Proper template-based preparation of one test protein gave
554 \code{HD} atoms and a 410:9 split of \code{N} versus \code{NA}, against a naive
writer that produced zero donors.

\item \pitfall{Feeding a truncated pocket file to a template-based preparation.}
Interrupted residues fail template matching. Prepare the whole protein, then
restrict the \emph{search box} to the pocket.

\item \pitfall{Deleting metals and cofactors indiscriminately.} A catalytic zinc is
often the anchor point of the whole site.

\item \pitfall{Ignoring protonation.} Both a ligand carboxylate and a receptor
histidine change hydrogen-bonding character depending on state, and neither is
determined by the input file.

\item \pitfall{Reading the affinity as a measured $\Delta G$.} It is an empirical
score with modest experimental correlation.

\item \pitfall{Trusting a single run.} Vina is stochastic, and multi-threaded runs
are not bitwise reproducible even with a fixed seed. Repeat, or raise
\code{exhaustiveness}, before concluding anything from a small difference.

\item \pitfall{Judging a docking setup by whether it crashes.} Every protocol in the
reference table ran cleanly on all 30 complexes while ranging from 43\% to 60\%
accuracy. Only re-docking distinguishes them.

\end{enumerate}

% =====================================================================
\section{Where to go next}
% =====================================================================

\begin{itemize}[leftmargin=1.3em, itemsep=3pt]
  \item \textbf{Re-scoring.} Because Vina often generates a good pose but ranks it
        poorly, re-scoring its pose set with a learned model is one of the
        highest-value additions to a docking pipeline.
  \item \textbf{Flexible side chains.} Vina supports selected flexible receptor
        residues, at a substantial cost in search difficulty.
  \item \textbf{Better charges.} Only relevant if you move to AutoDock4 scoring or
        to free-energy methods: the quality ladder runs Gasteiger $\rightarrow$
        MMFF94 $\rightarrow$ AM1-BCC $\rightarrow$ RESP.
  \item \textbf{Ensemble docking.} Docking into several receptor conformations
        partially compensates for Vina's rigid receptor.
  \item \textbf{Virtual screening.} Vina scores are weak at ranking \emph{across}
        ligands, so screening campaigns normally re-score or re-rank the top slice
        with a more expensive method.
\end{itemize}

\vspace{1em}
\noindent\rule{\linewidth}{0.4pt}\\[4pt]
{\small Vina: Trott \& Olson, \emph{J.\ Comput.\ Chem.}\ 31:455 (2010).
Vina 1.2: Eberhardt \emph{et al.}, \emph{J.\ Chem.\ Inf.\ Model.}\ 61:3891 (2021).
Meeko and the AutoDock suite: \url{https://autodock-vina.readthedocs.io}.}

\end{document}
